% !TeX program = xelatex
% !BIB program = biber
% Overleaf主文档：main.tex；编译器：XeLaTeX。
\def\mulu{./}
\makeatletter
\edef\input@path{{\mulu muban/}{\mulu}}
\makeatother
\documentclass{lunwen}\edef\wenxianlujing{\mulu wenxian.bib}
\expandafter\addbibresource\expandafter{\wenxianlujing}
% 前置内容：qianzhi.tex；后置内容：houzhi.tex；格式：muban/geshi.tex；字体：muban/ziti.tex。
% 正文层级：section=章，subsection=节，subsubsection=条，paragraph=款。
\begin{document}
% 前置内容：在这里填写论文信息、中英文摘要和关键词。
% 由main.tex加载，不单独编译；排版设置仍在geshi.tex。
% ==================== 01 论文信息 ==================== 中文标题正标题不超过25字，正副标题合计不超过35字。
\newcommand{\thesisTitle}{在此填写中文论文题目}
\newcommand{\thesisTitleEn}{YOUR THESIS TITLE IN ENGLISH}
\newcommand{\thesisType}{学术学位论文} % 专硕改为：专业学位论文
\newcommand{\authorName}{填写姓名}
\newcommand{\authorNameEn}{Your Name}
\newcommand{\supervisor}{填写导师及职称}
\newcommand{\supervisorEn}{Prof. Your Supervisor}
\newcommand{\degreeName}{工学硕士}
\newcommand{\degreeNameEn}{Master of Engineering}
\newcommand{\specialty}{建筑学}
\newcommand{\specialtyEn}{Architecture}
\newcommand{\submissionDate}{2026年12月}
\newcommand{\defenseDate}{2026年12月}
\newcommand{\defenseDateEn}{December, 2026}
\newcommand{\classification}{待填写}
\newcommand{\udc}{待填写}
\newcommand{\schoolCode}{待填写}


% ==================== 02 封面与内封 ====================
\makecover\makeinnertitles\frontmatter
% ==================== 03 中英文摘要 ====================
\begin{cnabstract}
在此填写中文摘要。依次概括研究背景、研究问题、研究内容与方法、主要结果和结论。摘要应独立完整，与英文摘要对应。
\keywords{关键词一；关键词二；关键词三}
\end{cnabstract}
\begin{enabstract}
Write the English abstract here. Its content should correspond to the Chinese abstract.
\enkeywords{keyword one, keyword two, keyword three}
\end{enabstract}

\clearpage\tableofcontents


\mainmatter

% ==================== 正文：在这里连续写作 ====================
\section{绪\quad 论}
\subsection{研究背景}
这里开始写你的正文。中文使用宋体，英文和数字使用 Times New Roman。正文自动首行缩进两个汉字，标题与图表自动编号。

这里演示顺序编码的上标引用\cite{lynch1960}。直接在句中说明文献时\cite{lynch1960}，可写“由文献\inlinecite{alexander1977}可知”。示例文献用于演示排版，请替换为你的实际文献。
\subsection{研究目的与意义}
\subsubsection{研究目的}
在此说明研究目的。
\subsubsection{研究意义}
在此说明研究意义。
\subsection{研究内容与方法}
本研究的分析流程如\figref{fig:framework}所示。
\begin{figure}[htbp]
 \centering
 \includegraphics[width=.8\linewidth]{\mulu tupian/example-liucheng.pdf}
 \figurenote{注：此图为排版示意，资料来源为作者绘制}
 \caption{研究流程示意}\label{fig:framework}
\end{figure}


\subsection{本章小结}
在此概括本章主要内容，小结中不放图表和公式。

\section{研究方法}
\subsection{研究数据}
数据字段如\tabref{tab:fields}所示。空缺值使用横线，表格采用三线表。
\begin{table}[htbp]
 \centering\tablefont
 \caption{数据字段及含义}\label{tab:fields}
 \begin{tabularx}{\linewidth}{lXr}
 \toprule 字段&含义&示例值\\\midrule
 ID&研究单元编号&001\\
 Area&研究单元面积（$\mathrm{m}^2$）&1200\\
 Count&观测数量&—\\\bottomrule
 \end{tabularx}
\end{table}
\subsection{模型构建}
\subsubsection{指标计算}
指标的计算方法见\eqnref{eq:performance}。
\begin{equation}
 P_i=\frac{N_i}{A_i}\qquad i=1,2,\ldots,10
 \label{eq:performance}
\end{equation}
\noindent\hangindent=3em\hangafter=1 式中\quad $P_i$——第 $i$ 个研究单元的指标值；\\
\hspace*{3em}$N_i$——该单元的观测数量；\\
\hspace*{3em}$A_i$——该单元面积，单位为 $\mathrm{m}^2$。
\paragraph{参数设定}
这里演示第四级标题，目录只列到第三级标题。
\subsection{跨页表格}
若表格超过一页，使用 longtable。示例见\tabref{tab:long}；续页自动重复编号和表头。
\begingroup\tablefont
\begin{longtable}{p{25mm}p{108mm}}
\caption{研究变量说明}\label{tab:long}\\
\toprule 变量&说明\\\midrule\endfirsthead
\noalign{\vskip 8bp}
\multicolumn{2}{r}{表\thetable（续表）}\\
\toprule 变量&说明\\\midrule\endhead
\bottomrule\endfoot
面积&研究单元的空间尺度指标\\
密度&研究单元单位面积的观测数量\\
可达性&研究单元与周边空间之间的连接程度\\
混合度&研究单元内不同功能的组合程度\\
\end{longtable}
\endgroup
\subsection{本章小结}
在此概括方法和数据。

% ==================== 06 结论 ====================
\conclusiontitle
在此填写结论。结论不加章序号，不标注文献引用，一般控制在2000字以内。
\begin{conclusions}
\item \textbf{在此填写第一项主要结论的标题}

在此展开说明第一项主要结论，正文使用宋体常规字重。
\item \textbf{在此填写第二项主要结论的标题}

在此展开说明第二项主要结论。
\end{conclusions}

% 后置内容：参考文献、创新成果、声明和致谢。
% 由main.tex加载，不单独编译；排版设置仍在geshi.tex。
\printbibliography[heading=bibintoc,title={参考文献}]
% ==================== 07 自动参考文献 ====================
% ==================== 08 创新成果 ====================
\achievementtitle
\achievementcategory{（一）发表的学术论文}
\begin{achievementlist}
\item 作者. 论文题目[J]. 期刊名称, 年, 卷(期): 起止页码.（对应本论文第1章）
\end{achievementlist}
\achievementcategory{（二）申请及已获得的专利}
\begin{achievementlist}
\item 发明人. 专利名称: 国别, 专利号[P]. 公开日期.
\end{achievementlist}
\achievementcategory{（三）参与的科研项目及获奖情况}
\begin{achievementlist}
\item 项目参与人. 项目名称, 项目来源, 课题编号.
\end{achievementlist}
% 没有的成果类别连同其列表一起删除；成果不得混入正文参考文献。

% ==================== 09 原创性声明与使用权限 ====================
\begin{thesisstatement}
\statementheading{学位论文原创性声明}
本人郑重声明：此处所提交的学位论文《\thesisTitle》，是本人在导师指导下，在哈尔滨工业大学攻读学位期间独立进行学术研究工作或专业实践工作所取得的成果，且学位论文中除已标注引用文献、资料的部分外不包含他人完成或已发表的成果；对本学位论文的学术研究工作或专业实践工作做出重要贡献的个人和集体，均已在文中以明确方式注明；对使用AI工具的情况均已在文中以明确方式标注。

\originalitysignature
\permissionheading
学位论文是研究生在哈尔滨工业大学攻读学位期间完成的成果，知识产权归属哈尔滨工业大学。学位论文的使用权限如下：
\begin{permissionclauses}
\item 学校可以采用影印、缩印或其他复制手段保存研究生上交的学位论文，并向国家图书馆报送学位论文；
\item 学校可以将学位论文部分或全部内容编入有关数据库进行检索和提供相应阅览服务；
\item 研究生毕业后发表与此学位论文研究成果相关的学术论文和其他成果时，应征得导师同意，且第一署名单位为哈尔滨工业大学。
\end{permissionclauses}
保密论文在保密期内遵守有关保密规定，解密后适用于此使用权限规定。

本人知悉学位论文的使用权限，并将遵守有关规定。

\permissionsignatures
\end{thesisstatement}

% ==================== 10 致谢 ====================
\section*{致\quad 谢}
在此填写致谢。


\end{document}
